\section{A Pin Is Three Numbers}
\label{sec:pin-three-numbers}

The type-state approach to GPIO is seductive: make an output pin a different \emph{type} from an
input pin, make a pin claimed by one driver unclaimable by another, and let the compiler refuse the
mistakes. This book's repository had exactly that — 605 lines of it — and it is now retired, for a
reason worth stating precisely.

\subsection{What the Small Version Looks Like}

The whole abstraction is a port, a bit and a polarity:

\begin{minted}{rust}
pub enum Port { A, B, C, D }

pub struct Output<'a> {
    gpio: &'a pac::gpio_s::RegisterBlock,
    port: Port,
    bit: u8,
    active_low: bool,
}

impl<'a> Output<'a> {
    pub fn on(&self)  { self.write(!self.active_low) }
    pub fn off(&self) { self.write(self.active_low) }
    pub fn toggle(&self) {
        let current = read_dout(self.gpio, self.port);
        self.write(current & (1 << self.bit) == 0);
    }
}
\end{minted}

Everything else is a register write:

\begin{minted}{rust}
fn write_dout(gpio: &pac::gpio_s::RegisterBlock, port: Port, bit: u8, high: bool) {
    // Read-modify-write, always: DOUT is one register for sixteen pins.
    let set = |r: u32, w: &mut pac::W<_>| {
        let mask = 1u32 << bit;
        let value = if high { r | mask } else { r & !mask };
        unsafe { w.bits(value) }
    };
    match port {
        Port::A => gpio.porta_dout().modify(set),
        Port::B => gpio.portb_dout().modify(set),
        Port::C => gpio.portc_dout().modify(set),
        Port::D => gpio.portd_dout().modify(set),
    }
}
\end{minted}

That is the entire output path: 40 lines including the four-way match. The input path is shorter
still — \codeinline{DIN} needs no mask, only a shift — and the mode setter is the read-modify-write
of \codeinline{MODEL} or \codeinline{MODEH} shown in the previous section.

\subsection{Why the Larger Version Was Retired}

The type-state layer was not wrong in principle. It was wrong here, in three ways that only show up
in a real project:

\begin{enumerate}
    \item \textbf{Nothing used it.} Its only consumer was a board module that a newer facade
          replaced. An API with zero users compiles, lints clean, and rots: at the time it was
          retired, its README documented the user LED as \codeinline{PB2} active-\emph{high}, when
          the board's LED is \codeinline{PA7} active-\emph{low}.
    \item \textbf{It moved the wiring into the HAL.} To hand out a typed \codeinline{LED} the HAL had
          to know which pin the LED was on — a board fact in the wrong layer. The layered rule from
          section~\ref{sec:layer-boundaries} is that the HAL names no board wiring, and the typed
          pins were the mechanism by which it did.
    \item \textbf{It cost more than it prevented.} The three real defects in this project's GPIO
          code — a blind \codeinline{MODEL} write, a mode value with the wrong width, and an
          open-drain pin configured push-pull — were all invisible to the type system, because all
          three are legal code with the wrong \emph{number} in it. Types catch the mistakes you can
          name; registers catch the rest.
\end{enumerate}

\subsection{What Replaced the Type System}

An ownership rule that is checkable without generics: the board takes the peripherals once
(\codeinline{Board::take()}), and every pin comes out of it with a port and a bit named at the call
site:

\begin{minted}{rust}
let board = Board::take().expect("peripherals taken twice");
let led     = board.output_on(Port::A, 7, true);      // active low
let button  = board.input_on(Port::C, 1);             // pull-up fitted on the base
let buzzer  = board.output_on(Port::C, 3, false);
\end{minted}

The exclusivity that the type system used to provide is now the PAC singleton's job — one
\codeinline{take()} per program — and the wiring that the HAL used to hide lives in whichever crate
owns the board. Fewer lines, fewer layers, and the two facts a reader needs (which pin, which
polarity) are in the code they are reading.
